\documentclass[12pt]{article}
\usepackage[utf8]{inputenc}
\usepackage[T1]{fontenc}
\usepackage{hyperref}
\title{Wares Language Resources for K-12 Curriculum Development}
\author{Compiled by Claude Sonnet 5 (claude-sonnet-5) for Andr\'as Kornai}
\date{}
\begin{document}
\maketitle

\section{Introduction}
Wares (ISO 639-3 wai; Glottocode ware1254) is an endangered language of Indonesia, reported to belong to the Tor-Kwerba family. Searches for this language returned no page about Wares itself. The results covered only neighbouring Tor-Kwerba languages (Kwesten, Kwinsu, Kwerba, Kauwera, Isirawa), and none of them is Wares. Nothing below should be read as a Wares resource.

\section{Primary Educational Resources}
None found. No primers, textbooks or school programmes for Wares appeared in search results.

\section{Linguistic and Reference Materials}
None found. No Wares dictionary, grammar or orthography was located.

\section{Religious and Community Resources}
None found.

\section{Digital and Audio Resources}
None found.

\section{Curriculum Development Context}
No curriculum can be grounded in the sources found. Further work would need fieldwork records, Ethnologue or Glottolog bibliographies, or archives such as Paradisec or ELAR, none of which were checked in this search. The language is thought to be spoken as a first language by adults only, so documentation would have to come before any school use.
\end{document}